\section{Introduction}
\label{sec:intro}

Public leaderboards turn a score matrix of models by tasks into a single ranking, and the top place carries visibility.
The ranking depends on the aggregation rule, and developers can raise a model's score on chosen tasks, by targeted data or tuning.
We ask a counterfactual question about the scores: \emph{how much targeted improvement does a given model need to reach first place, and how does that depend on the rule that aggregates the tasks?}
We answer it for a single target model with all rivals held fixed, on two public matrices, the Open LLM Leaderboard (OLLB; 4{,}496 models, six tasks) and HELM MMLU (79 models, 57 subjects).

The question has three parts that earlier work on task selection and aggregation does not separate (Section~\ref{sec:related}).
First, ``how much'' is not one number: we report the total points added ($L_1^\ast$), the number of tasks touched ($K^\ast$) and the largest single-task lift ($L_\infty^\ast$) as three separate optima.
Second, an investment is meaningful only under a cap on how far each score can rise, and under attainable caps first place is often out of reach; we therefore report \emph{feasibility} (whether first place is attainable at all, and the smallest cap or uniform lift that makes it so) as the headline, and compare costs only on the cohort of targets that are feasible under every compared rule.
Third, the answer may depend on the score scale and the number of tasks, which we vary.

\paragraph{Contributions.}
(i) We formalise the score-edit counterfactual and its three optima, with feasibility kept separate from cost, for score rules (mean, median, minimum, geometric mean) and ordinal rules (Borda, strict Condorcet winner; the latter is an upper bound for any Condorcet-consistent rule).
(ii) We give closed forms for the mean, median (odd $n$) and minimum, and show that $L_\infty^\ast$ equals the smallest uniform lift for every rule; we check them against exact solvers, including integer programs for the ordinal rules, on 31{,}080 instances, all of which agree (Section~\ref{sec:method}, Appendix~\ref{app:proofs}).
(iii) We run all rules over eight cap scenarios on 4{,}495 OLLB and 78 HELM targets (Section~\ref{sec:exp}).

\paragraph{Findings.}
They are conditional on the cap and the scale and are stated as such.
\textbf{(F1)} If no task may be raised by more than 10 points, first place is attainable for only 3--29\% of OLLB targets and 18--51\% of HELM targets, depending on the rule; for models ranked 2--30 under the mean the shares are 76--100\% and 41--100\%.
Rules differ in \emph{how many} models can win, not only in what it costs: on OLLB, Borda reaches 29\% at 10 points where the mean reaches 3\%.
\textbf{(F2)} The ordering of rules is not stable across datasets and caps: the minimum rule is the hardest to satisfy under small caps on HELM (18\% at 10 points) and the easiest on OLLB under a 20-point cap (64\% against 14--21\% for the mean, median and geometric mean).
\textbf{(F3)} Under loose caps, on the cohort of targets feasible under every rule, the minimum rule needs the fewest total points on OLLB and, with the median, on HELM, while Borda needs the most tasks (5 of 6 on OLLB, 46 of 57 on HELM).
\textbf{(F4)} On OLLB, whose six tasks have unrelated scales, the number of tasks needed by targets ranked beyond 10 changes by up to four with the score scale (official, raw, $z$-score, rank percentile), so OLLB conclusions are tied to the official scale.
\textbf{(F5)} The closed forms explain these differences: the minimum rule requires every task below the best rival's level to reach it, so it is blocked by any binding cap there.

\emph{Investment here is a score counterfactual, not a training cost}; it ignores spillover between tasks, rival responses and the data or compute that would deliver the gain (Section~\ref{sec:limits}).
